\documentclass[aspectratio=169,12pt,t]{beamer}
\usepackage{graphicx}
\setbeameroption{hide notes}
\setbeamertemplate{note page}[plain]
\usepackage{listings}
\usepackage{eepic}

% header.tex: boring LaTeX/Beamer details + macros

% get rid of junk
\usetheme{default}
\beamertemplatenavigationsymbolsempty
\hypersetup{pdfpagemode=UseNone} % don't show bookmarks on initial view
\usepackage[export]{adjustbox}


% font
\usepackage{fontspec}
\setsansfont
  [ ExternalLocation = fonts/ ,
    UprightFont = *-regular ,
    BoldFont = *-bold ,
    ItalicFont = *-italic ,
    BoldItalicFont = *-bolditalic ]{texgyreheros}
\setbeamerfont{note page}{family*=pplx,size=\footnotesize} % Palatino for notes
% "TeX Gyre Heros can be used as a replacement for Helvetica"
% I've placed them in fonts/; alternatively you can install them
% permanently on your system as follows:
%     Download http://www.gust.org.pl/projects/e-foundry/tex-gyre/heros/qhv2.004otf.zip
%     In Unix, unzip it into ~/.fonts
%     In Mac, unzip it, double-click the .otf files, and install using "FontBook"

% named colors
\definecolor{offwhite}{RGB}{255,250,240}
\definecolor{gray}{RGB}{155,155,155}
\definecolor{purple}{RGB}{177,13,201}
\definecolor{green}{RGB}{46,204,64}

\definecolor{background}{RGB}{255,255,255}
\definecolor{foreground}{RGB}{24,24,24}
\definecolor{title}{RGB}{27,94,134}
\definecolor{subtitle}{RGB}{22,175,124}
\definecolor{hilit}{RGB}{122,0,128}
\definecolor{vhilit}{RGB}{255,0,128}
\definecolor{codehilit}{RGB}{255,0,128}
\definecolor{lolit}{RGB}{95,95,95}
\definecolor{myyellow}{rgb}{1,1,0.7}
\definecolor{nhilit}{RGB}{128,0,128}  % hilit color in notes
\definecolor{nvhilit}{RGB}{255,0,128} % vhilit for notes

\newcommand{\hilit}{\color{hilit}}
\newcommand{\vhilit}{\color{vhilit}}
\newcommand{\nhilit}{\color{nhilit}}
\newcommand{\nvhilit}{\color{nvhilit}}
\newcommand{\lolit}{\color{lolit}}

% use those colors
\setbeamercolor{titlelike}{fg=title}
\setbeamercolor{subtitle}{fg=subtitle}
\setbeamercolor{institute}{fg=lolit}
\setbeamercolor{normal text}{fg=foreground,bg=background}
\setbeamercolor{item}{fg=foreground} % color of bullets
\setbeamercolor{subitem}{fg=lolit}
\setbeamercolor{itemize/enumerate subbody}{fg=lolit}
\setbeamertemplate{itemize subitem}{{\textendash}}
\setbeamerfont{itemize/enumerate subbody}{size=\footnotesize}
\setbeamerfont{itemize/enumerate subitem}{size=\footnotesize}

% page number
\setbeamertemplate{footline}{%
    \raisebox{5pt}{\makebox[\paperwidth]{\hfill\makebox[20pt]{\lolit
          \scriptsize\insertframenumber}}}\hspace*{5pt}}

% add a bit of space at the top of the notes page
\addtobeamertemplate{note page}{\setlength{\parskip}{12pt}}

% default link color
\hypersetup{colorlinks, urlcolor={hilit}}

\lstset{language=bash,
        basicstyle=\ttfamily\scriptsize,
        frame=single,
        commentstyle=,
        backgroundcolor=\color{offwhite},
        showspaces=false,
        showstringspaces=false
        }


% a few macros
\newcommand{\bi}{\begin{itemize}}
\newcommand{\bbi}{\vspace{24pt} \begin{itemize} \itemsep8pt}
\newcommand{\ei}{\end{itemize}}
\newcommand{\be}{\begin{enumerate}}
\newcommand{\bbe}{\vspace{24pt} \begin{enumerate} \itemsep8pt}
\newcommand{\ee}{\end{enumerate}}
\newcommand{\sbi}{\begin{itemize} \fontsize{9pt}{9.5}\selectfont}
\newcommand{\sbe}{\begin{enumerate} \fontsize{9pt}{9.5}\selectfont}
\newcommand{\ig}{\includegraphics}
\newcommand{\subt}[1]{{\footnotesize \color{subtitle} {#1}}}
\newcommand{\ttsm}{\tt \small}
\newcommand{\ttfn}{\tt \footnotesize}
\newcommand{\figh}[2]{\centerline{\includegraphics[height=#2\textheight]{#1}}}
\newcommand{\figw}[2]{\centerline{\includegraphics[width=#2\textwidth]{#1}}}
\newcommand{\fighboxed}[2]{\centerline{\includegraphics[frame,height=#2\textheight]{#1}}}
\newcommand{\figwboxed}[2]{\centerline{\includegraphics[frame,width=#2\textwidth]{#1}}}


%%%%%%%%%%%%%%%%%%%%%%%%%%%%%%%%%%%%%%%%%%%%%%%%%%%%%%%%%%%%%%%%%%%%%%
% end of header
%%%%%%%%%%%%%%%%%%%%%%%%%%%%%%%%%%%%%%%%%%%%%%%%%%%%%%%%%%%%%%%%%%%%%%

% title info
\title{data cleaning principles}
\author{\vspace*{-10pt} \href{https://kbroman.org}{Karl Broman}}
\institute{Biostatistics \& Medical Informatics, UW{\textendash}Madison}
\date{\href{https://twitter.com/kwbroman}{\tt \scriptsize \color{foreground} @kwbroman}
\\[-4pt]
\href{https://kbroman.org}{\tt \scriptsize \color{foreground} kbroman.org}
\\[-4pt]
\href{https://github.com/kbroman}{\tt \scriptsize \color{foreground} github.com/kbroman}
\\[-4pt]
{\scriptsize \href{https://kbroman.org/Talk_DataCleaning}{\tt kbroman.org/Talk\_DataCleaning}}
}


\begin{document}

% title slide
{
\setbeamertemplate{footline}{} % no page number here
\frame{
  \titlepage

  \vfill \hfill
  \href{https://creativecommons.org/publicdomain/zero/1.0/}{\includegraphics[height=7mm]{Figs/cc-zero.png}}

  \note{
    These are slides for a talk for the csv,conf,v6 ({\tt
    https://csvconf.com/}) on May 4--5, 2021.

    Data analysts spend a lot of time organizing and cleaning data,
    but few of us have been trained to do so. Why is that?

    Some say that data cleaning is difficult to generalize. But I
    think there are some general principles. Moreover, I think we have
    an important shared experience in data cleaning that we can
    commiserate about, and through which we can learn from each other.
  }

} }


\begin{frame}[c]{}

\centering
\Large

Tidy data are all alike, \\[3pt]
but every messy dataset  \\[3pt]
is messy in its own way.

\bigskip\bigskip

\hspace{5cm} --
\href{https://r4ds.had.co.nz/tidy-data.html}{Hadley Wickham}



\note{
  Hadley's talking more about data organization than data
  cleanliness.  And his point is that if you make data tidy, it
  simplifies all the downstream analyses.

  But {\hilit is} every messy dataset {\vhilit uniquely} messy?

  For sure, many my collaborators have shown impressive creativity in
  their approach to organizing and managing data. But we do see
  many of the same sorts of problems over and over.
}
\end{frame}


\begin{frame}[c]{}

\centering
\Large

If I clean up [Medicare] data ...\\[3pt]
does any of the knowledge I gain ...\\[3pt]
apply to the processing of RNA-seq data?


\bigskip\bigskip

\hspace{5cm} --
\href{https://doi.org/10.1080/10618600.2017.1385470}{Roger Peng}



\note{
  In his discussion of David Donoho's paper about data science, Roger
  Peng wrote about how data cleaning is frustratingly difficult to
  generalize.

  But my answer to his question is {\vhilit absolutely!}

  A person with experience cleaning one dataset has important
  experience to draw upon when moving to another dataset, even if it's
  of a totally different nature.
}
\end{frame}



\begin{frame}[c]{}


  \href{https://datamishapsnight.com}{\figh{Figs/data_mishaps_night.png}{1.0}}

  \vspace{-7mm}
  \hfill
  \begin{minipage}[c]{0.33\textwidth}
    {\footnotesize Caitlin Hudon \& Laura Ellis}

    \href{https://datamishapsnight.com}{\lolit \tt dataMishapsNight.com}
  \end{minipage}

\note{
  In February, 2021, Caitlin Hudon and Laura Ellis organized a
  Friday evening conference where 16 people gave short presentations on data
  mishaps.

  Many of the stories concerned mistakes in data cleaning, and these
  seemed to bring out a strong sense of shared experience. We have
  suffered and struggled through very similar data problems.
}
\end{frame}



\begin{frame}{Data cleaning}


\vspace*{-6pt}

\begin{minipage}[t]{0.45\textwidth}

\bbi
\item tedious
\item embarrassing
\item needs context
\item doesn’t feel like progress
  \ei

\end{minipage} \hfill
\begin{minipage}[t]{0.45\textwidth}

\only<2>{
\bbi
\item requires creativity
\item requires coding prowess
\item source of most problems
\ei
}
\end{minipage}

\note{
  Really, I think we don't usually teach data cleaning
  because it's something we prefer to keep private.

  We're shy about it.

  And data cleaning code is our ugliest code.
}
\end{frame}



\begin{frame}[c]{}

  \figh{Figs/logo.pdf}{0.8}

\note{
  I'm proposing a set of basic principles for data cleaning, and
  splitting them into five groups. There are some fundamental
  principles, followed by four basic ideas: verify things that you
  expect, explore to find further oddities, ask questions, and
  document what you've done.
}
\end{frame}


\begin{frame}{}
\lfigh{Figs/logo_fundamentals.pdf}{0.2}

\bigskip\bigskip\bigskip

\centering
\Large

{\princolorA 1. Don't clean data when you're tired or hungry.}

\bigskip\bigskip

\large
\hspace{5cm}
(paraphrasing \href{https://twitter.com/ghazalgulati}{Ghazal Gulati})


\note{
  At her talk at the Data Mishaps night, Ghazal Gulati emphasized this
  point, of not cleaning data when you're tired or hungry.

  Data cleaning requires considerable concentration, and you need to
  allow sufficient time to do the work. If you're in a hurry, you'll
  miss things.
}
\end{frame}


\begin{frame}{}
\lfigh{Figs/logo_fundamentals.pdf}{0.2}

\bigskip\bigskip\bigskip


\only<1|handout 0>{
\centering
\Large

{\princolorA 2. Don't trust anyone (even yourself)}
}

\only<2>{
  \vspace*{-27.5mm}
  \hspace{0.4\textwidth} {\princolorA 2. Don't trust anyone (even yourself)}

\vspace{0.25\textheight}

\centering
\large

``my motto is `trust no one' \hspace{1cm} \\[3pt]
\hspace{1cm} ...except maybe @kwbroman?''

\bigskip

\hspace{5cm}
-- \href{https://twitter.com/JennyBryan/status/595689693162938368}{Jenny Bryan}
}

\note{
   Next: don't trust anyone. Even if the initial data cleaning was
   done by someone you respect, you should double-check things that
   they may have missed. And data cleaning is an ongoing process.

   Jenny Bryan's tweet is among the nicest things anyone has said
   about me.
}
\end{frame}



\begin{frame}{}
\lfigh{Figs/logo_fundamentals.pdf}{0.2}

  \vspace*{-15mm}
  \hspace{0.4\textwidth} {\princolorA 3. Think about what might have gone wrong} \\
  \hspace{0.4\textwidth} {\color{background} 3.}
            {\princolorA and how it might be revealed} \hspace{20.5mm}

\vspace{0.06\textheight}

\figh{Figs/sample_mixups.png}{0.78}



\note{
  Personally, I think this is the most important principle for data
  cleaning. It has been central in guiding my approach.

  The figure here is an illustration of the most startingly result
  I've had in data cleaning: a genetics project where 25\% of the
  DNA samples had been mislabelled. The samples were arranged in wells
  in 8$\times$12 plates; four of the six plates are shown here. The
  dots indicate the correct DNA was placed in the well, but the arrows
  point from where a sample should have been to where it actually was
  placed.

  I ultimately came to this finding by thinking about what might have
  gone wrong in the project, checking for particularly problems, and
  then following the trail of evidence to this mess.
}
\end{frame}


\begin{frame}{}
\lfigh{Figs/logo_fundamentals.pdf}{0.2}

\vspace{-15mm} \hspace{0.4\textwidth}
{\princolorA 4. Use care in merging}


\vspace{0.12\textheight}

\includegraphics[height=0.6\textheight]{Figs/spreadsheet_colnames.pdf}

\vspace*{-0.45\textheight}
\hspace*{0.22\textwidth}
\includegraphics[height=0.6\textheight]{Figs/spreadsheet_colnames2.pdf}


\note{
  Many problems arise due to mistakes when merging data from multiple
  files. A common problem is a change in the data arrangement, such as
  in the order of columns.

  Focus on the labels (which are more likely correct), rather than the
  position of variables in a file (which are more likely to change).
}
\end{frame}


\begin{frame}{}
\lfigh{Figs/logo_fundamentals.pdf}{0.2}

\bigskip\bigskip\bigskip

\centering
\Large

{\princolorA 5. Dates \& categories suck}


\note{
  The fifth fundamental principle is that dates and categories suck.
  You'll expend an inordinate amount of time dealing with these: typos
  in category labels, different date formats, people who died a decade
  they were born or lived to be 150.

  Just be glad if you're not dealing with time zones.

  But you may be asking yourself, ``How is this a principle.''

}
\end{frame}


\begin{frame}{}

\bigskip \bigskip \bigskip \bigskip

\Large

{\color{title} Principle:} \hspace{50mm}

\centering

\bigskip \bigskip

{\hilit a fundamental truth that guides our thinking}


\note{
  Yeah, are these principles? I was thinking the same thing. Was I
  drifting away from principles and more to just stuff to know or do?

  This seems a pretty good definition, and is sufficiently broad to
  cover what I'm proposing, for the most part.
}
\end{frame}


\begin{frame}{}
\lfigh{Figs/logo_fundamentals.pdf}{0.2}

\bigskip\bigskip\bigskip

\centering
\Large

{\princolorA 5. Dates \& categories suck}


\note{
  So yeah, this counts as a principle.

  Much of your pain will come from the dates and categorical data;
  you should be ready for that.
}
\end{frame}




\begin{frame}{}
\lfigh{Figs/logo_verify.pdf}{0.2}

\vspace{-15mm} \hspace{0.4\textwidth}
{\princolorB 6. Check that distinct things are distinct}

\vspace{0.16\textheight}

\figh{Figs/dup_ids.pdf}{0.7}


\note{
   The nature major section concerns that you should go through and
   {\nhilit verify} all of the things that you know should be true.

   First up: are the variables that are supposed to have distinct
   values really showing distinct values? Here, there are a pair of
   individual identifiers that are duplicated.
}
\end{frame}


\begin{frame}{}
\lfigh{Figs/logo_verify.pdf}{0.2}

\vspace{-15mm} \hspace{0.4\textwidth}
{\princolorB 7. Check that matching things match}

\vspace{0.22\textheight}

\begin{minipage}[t]{0.48\textwidth}

\includegraphics[width=\textwidth]{Figs/mismatchB.pdf}
\end{minipage}
\hfill
\begin{minipage}[t]{0.48\textwidth}

\includegraphics[width=\textwidth]{Figs/mismatch.pdf}
\end{minipage}


\note{
    Next, are the things that are supposed to match actually matching?
    This often concerns data that are repeated in multiple data files.
    For example, you might have body weight included multiple times;
    are individuals' weights the same in all files? Here, there's a
    column for ``number of generations'' that shows an inconsistency
    between two files.
}
\end{frame}



\begin{frame}{}
\lfigh{Figs/logo_verify.pdf}{0.2}

\vspace{-15mm} \hspace{0.4\textwidth}
{\princolorB 8. Check calculations}

\vspace{0.18\textheight}

\figh{Figs/check_calculations.pdf}{0.72}

\note{
    Any time there is a calculation, you should verify the values.
    This is useful both for ensuring that you understand the
    calculation, and you also have the chance to reduce round-off
    error (though that seldom matters). Occasionally you may see
    mistakes.

    Here, I'm plotting HOMA-IR which is the product of blood glucose
    and blood insulin; the values in the data on the x-axis and the
    values I calculated on the y-axis. It's often better to look at the
    differences directly; in the right panel I plot the differences on
    the y-axis. It appears that there is a batch of individuals that
    whose values were rounded more coarsely.
}
\end{frame}



\begin{frame}{}
\lfigh{Figs/logo_verify.pdf}{0.2}

\centering
\Large

\bigskip \bigskip \bigskip

{\princolorB 9. Look for other instances of a problem}

\note{
    As with software testing, any time you find a problem, be sure to
    look for other instances of that problem.
}
\end{frame}



\begin{frame}{}
\lfigh{Figs/logo_explore.pdf}{0.2}

\vspace{-15mm} \hspace{0.4\textwidth}
{\princolorC 10. Make lots of plots}

\vspace{0.15\textheight}

\only<1>{\figh{Figs/il3.pdf}{0.75}}
\only<2|handout 0>{\figh{Figs/body_weight.pdf}{0.75}}
\only<3|handout 0>{\figh{Figs/adipose_weight.pdf}{0.75}}
\only<4|handout 0>{\figh{Figs/spreadsheet_g_v_mg.pdf}{0.75}}


\note{
}
\end{frame}


\begin{frame}{}
\lfigh{Figs/logo_explore.pdf}{0.2}

\vspace{-15mm} \hspace{0.4\textwidth}
{\princolorC 11. Look at missing value patterns}

\vspace{0.12\textheight}

\figh{Figs/scatter_na.pdf}{0.8}

\note{
visdat: https://docs.ropensci.org/visdat/

naniar: http://naniar.njtierney.com/
}
\end{frame}


\begin{frame}{}
\lfigh{Figs/logo_explore.pdf}{0.2}

\vspace{-15mm} \hspace{0.4\textwidth}
       {\princolorC 12. With massive data,} \\
\hspace{0.4\textwidth} {\color{background} 12.} {\princolorC make more plots not fewer}

\vspace{0.10\textheight}

\only<1|handout 0>{\figh{Figs/hypo_arrays.pdf}{0.75}}
\only<2|handout 0>{\figh{Figs/hypo_arrays_hilit.pdf}{0.75}}
\only<3|handout 0>{\figh{Figs/hypo_iqr_v_median.pdf}{0.75}}
\only<4>{\figh{Figs/hypo_boxplots.pdf}{0.75}}


\note{
  With large-scale datasets, it can be hard to make the sort of
  exploratory plots that you'd typically make. With oodles of data,
  you'd think you'd be looking at oodles of plots, but there's a
  tendency to give up and not look at any.

  It's hard to look at 500 histograms, but it can be done. Superimpose
  a bunch of density estimates, maybe highlighting some portion of
  them. You can also pull out a couple of summary statistics, such as
  the median and inter-quartile range.

  Or here I'm looking at the equivalent of 500 boxplots. I sorted a
  set of gene expression microarrays by their median, and then plotted
  the median in blue, the 25th and 75th percentile in black, the 10th
  and 90th in green, 5th and 95th in red, and 1st and 99th in purple.

  With these data, it became apparent that there were 120 badly
  behaved arrays, with median shifted to the right and with a long
  left tail.
}
\end{frame}


\begin{frame}{}
\lfigh{Figs/logo_explore.pdf}{0.2}

\vspace{-15mm} \hspace{0.4\textwidth}
{\princolorC 13. Follow up all artifacts}


\vspace{0.10\textheight}

\figh{Figs/weird_correlation_matrix.png}{0.75}

\hfill
{
\scriptsize
\href{https://kbroman.org/blog/2012/04/25/microarrays-suck/}{\tt kbroman.org/blog/2012/04/25/microarrays-suck}
}

\note{
  Wow the clash of those colors is particularly bad.

  This is a heat map of the correlation matrix for a set of gene
  expression microarrays. The plaid pattern was a shock to me, and
  was caused by a set of bad arrays that we hadn't noticed previously.

  My point here is simply to follow up all artifacts.

  If you see something weird, follow through and try to figure out the
  underlying cause. If could be an error, or a set of bad assays, or
  it could be the most interesting finding in the study.
}
\end{frame}


\begin{frame}{}
\lfigh{Figs/logo_ask.pdf}{0.2}

\hfill
\begin{minipage}[t]{0.9\textwidth}


\bbi
\princolorD
\item[\princolorD 14.] Ask questions
\item[\princolorD 15.] Ask for the primary data
\item[\princolorD 16.] Ask for metadata
\item[\princolorD 17.] Ask why data are missing
\ei

\end{minipage}

\note{
}
\end{frame}


\begin{frame}{}
\lfigh{Figs/logo_document.pdf}{0.2}

\hfill
\begin{minipage}[t]{0.9\textwidth}


\bbi
\princolorE
\item[\princolorE 18.] Create checklists \& pipelines
\item[\princolorE 19.] Document not just what but why
\item[\princolorE 20.] Expect to recheck
\ei

\end{minipage}

\note{
}
\end{frame}



\begin{frame}[c]{}

  \figh{Figs/logo_withtext.pdf}{0.95}

\note{
  In summary...
}
\end{frame}



\begin{frame}[c]{}

\Large

Slides: \href{https://kbroman.org/Talk_DataCleaning}{\tt kbroman.org/Talk\_DataCleaning}

\vspace*{-7mm}
\hfill
\href{https://creativecommons.org/publicdomain/zero/1.0/}{\includegraphics[height=7mm]{Figs/cc-zero.png}}

\vspace{5mm}

\href{https://kbroman.org}{\tt \lolit kbroman.org}

\vspace{7mm}

\href{https://github.com/kbroman}{\tt \lolit github.com/kbroman}

\vspace{7mm}

\href{https://twitter.com/kwbroman}{\tt \lolit @kwbroman}

\note{
}

\end{frame}




\end{document}
